% !TeX root = ../main.tex
\section{Regeneration and Pattern Formation}
\label{sec:regeneration-pattern-formation}

% TODO: add biological references for wound-induced regeneration.
% TODO: add references for observed activator/Wnt peaks.
% TODO: add numerical figures for the scenarios described below.
% TODO: specify the final domain, boundary conditions, and numerical protocol.

We now ask what can already be explained by the two-equation reference model
\eqref{eq:gm-inhibitor-production}, before introducing source density. This is
the minimal test that the model should pass: if basic regeneration and peak
formation already follow from the activator-inhibitor mechanism, then source
density should not be credited with explaining these phenomena. Its role must be
looked for elsewhere.

In the biological interpretation, the activator \(u\) represents a local
head-forming signal. A regenerated head is therefore represented, at this level
of description, by the emergence of a localized activator peak. The reference
model has exactly the phase-space structure needed for a threshold response: the
zero state is stable, a lower positive state acts as a threshold, and an upper
positive state is the branch from which spatial patterns can form. This makes it
possible to distinguish weak perturbations, which decay, from sufficiently
strong wound-like perturbations, which trigger regeneration.

If the system is initialized close to the diffusion-unstable homogeneous state,
then small spatial perturbations grow and a pattern is formed. This is the
standard pattern-formation regime. It shows that the model can create localized
activator peaks, but it does not yet represent a regeneration experiment,
because the system is already placed near a pattern-forming state.

A regeneration experiment is better represented by starting near the resting
state. If the head is removed but no additional wound signal is introduced, the
perturbation is too weak to cross the threshold. The solution then returns to the
zero steady state, and no new head is formed. This is a useful feature rather
than a failure: it means that the model does not regenerate automatically from
every small disturbance.

Regeneration appears when the removal experiment is accompanied by a sufficiently
strong localized perturbation. Such a perturbation represents a wound signal. If
it pushes the solution past the unstable threshold branch, an activator peak can
grow and persist. In this sense, the two-equation model already captures the
basic logic of wound-induced head formation: injury is not merely a change in
geometry, but a local activation event strong enough to move the tissue out of
the basin of attraction of the resting state.

The same mechanism can also produce two-head configurations on sufficiently
large domains. A second localized peak can be initiated by introducing a wound in
the middle of the tissue or near one of the ends. However, the final position of
the new peak is not arbitrary. The model selects preferred patterning positions,
for example at a boundary, near the middle, or close to one third of the tissue,
depending on the domain size and on the unstable spatial modes available to the
system. A wound imposed away from these preferred locations may create a
transient peak at the wound site, but this peak can later drift toward a
model-selected position.

This behaviour is especially visible when the domain size is varied. On a small
domain, simulations often converge to a single boundary peak. If this one-peak
solution is continued while the domain is gradually increased, it can remain
stable even on much larger domains. By contrast, if a large domain is initialized
near the zero state and a strong boundary perturbation is introduced, the system
again tends to form a localized boundary peak. But if the same large domain is
initialized close to the diffusion-unstable homogeneous state, a small
perturbation typically does not select a single head. Instead, the solution
converges to a multi-peak pattern.

Thus the observed number and position of peaks are not determined only by the
parameters of the equations. They also depend on the domain size, the initial
condition, and the path by which the pattern is generated. A one-head pattern
obtained by continuation from a small domain is not the same numerical experiment
as a pattern selected from small random perturbations of a homogeneous
pattern-forming state.

The conclusion is that source density is not needed for these basic regeneration
phenomena. The two-equation Gierer--Meinhardt model already accounts for decay
of weak perturbations, wound-induced head formation, one- and two-peak patterns,
and the selection of preferred peak locations. What it does not provide is an
autonomous mechanism for positional memory: the model can create peaks, but it
does not remember where a tissue fragment came from, nor does it guarantee that a
head remains anchored at an arbitrary graft-induced position. This is the point
at which source density becomes relevant.
